\begin{definition}[Divisible abelian group] \label{definition:divisible_commutative_semigroup_divisible_commutative_monoid_divisible_abelian_group}
    A \CrefAndHyperrefIfExist{definition:monoid}{commutative semigroup} $S$, say written additively, is called \hldef{divisible} if for every $s \in S$ and every positive integer $n \in \bbZ_{>0}$, there exists some $s' \in S$ such that $s' = ns$\CrefIfExists{definition:power_of_a_semigroup_element_by_a_positive_integer}.
    % A \hldef{divisible abelian group} is an \CrefAndHyperrefIfExist{definition:group}{abelian group} $D$ such that for every $d \in D$ and every nonzero integer $n$, there exists $x \in D$ satisfying
    % $$nx = d.$$
    Equivalently, the multiplication-by-$n$ map
    $$[n] : S \longrightarrow S, \quad s \longmapsto ns$$
    is \CrefAndHyperrefIfExist{definition:surjective_function}{surjective} for every positive integer $n$.

    In particular, we may also speak of \hldef{divisible commutative monoids}, and \hldef{divisible commutative groups}

    % A commutative monoid $M$ is called a \hldef{divisible monoid} if the same condition holds with the abelian group replaced by a monoid.
\end{definition}
